\documentclass[11pt]{article}
\usepackage[margin=1in]{geometry}
\usepackage{amsmath,amssymb,amsthm}
\usepackage{booktabs}
\usepackage[hidelinks]{hyperref}

\newtheorem{theorem}{Theorem}
\newtheorem{lemma}[theorem]{Lemma}
\newtheorem{corollary}[theorem]{Corollary}

\newcommand{\R}{\mathbb{R}}
\newcommand{\C}{\mathbb{C}}
\newcommand{\PP}{\mathbb{P}}
\newcommand{\cM}{\mathcal{M}}
\newcommand{\cN}{\mathcal{N}}
\newcommand{\cF}{\mathcal{F}}
\newcommand{\cL}{\mathcal{L}}
\newcommand{\cD}{\mathcal{D}}
\newcommand{\cH}{\mathcal{H}}
\newcommand{\cX}{\mathcal{X}}
\DeclareMathOperator{\mult}{mult}

\title{Covering numbers of semi-algebraic sets:\\
linear in the dimension, logarithmic in the degree}
\author{}
\date{}

\begin{document}
\maketitle

\begin{abstract}
Let $\cM\subseteq\R^N$ be a semi-algebraic set of dimension $m$ whose Zariski
closure has degree $d$. For every Euclidean norm we show
$\log\cN_\varepsilon(\cM\cap B)\le\log d+2m\log(1+2/\varepsilon)$, with no
dependence on $N$. The proof applies Gromov's packing argument for the entropy
of holomorphic maps to the complexification of $\cM$, using Lelong's
monotonicity theorem. Since evaluation at data points is linear, the bound also
holds for empirical covering numbers; for polynomial machine learning models
this gives a sample complexity of order
$\varepsilon^{-2}\bigl(m+\log d+\log(1/\delta)\bigr)$.
\end{abstract}

\section{Setting and results}

Let $V$ be a real vector space of dimension $N$, for instance a space of
polynomials of bounded degree containing the neuromanifold of a polynomial
model~\cite{MSMTK25}, and let $\cM\subseteq V$ be semi-algebraic of dimension
$m$. Let $X\subseteq V_\C=V\otimes_\R\C$ be the complex Zariski closure of
$\cM$, with irreducible components $X_1,\dots,X_s$. Then $\dim_\C X=m$, since
a semi-algebraic set has the dimension of its Zariski
closure~\cite[\S2.8]{BCR98} and Krull dimension is unchanged by extending
scalars to $\C$. We set
\[
  d:=\deg X:=\sum_{j=1}^{s}\deg X_j,
\]
where $\deg X_j$ is the number of points of $X_j\cap L$ for a generic affine
subspace $L$ of codimension $\dim X_j$. If $\cM=\varphi(\R^q)$ for a polynomial
map $\varphi$, as for polynomial neural networks, then $X$ is irreducible and
$d$ is the degree of the neurovariety; it bounds the number of complex
interpolants of $m$ generic samples~\cite[App.~B]{MSMTK25}. For a subset $S$ of
a normed space, $\cN_\varepsilon(S)$ is the least number of closed
$\varepsilon$-balls covering~$S$.

\begin{theorem}\label{thm:main}
Fix an inner product on $V$, with closed unit ball $B$. For every
$\varepsilon>0$,
\[
  \cN_\varepsilon(\cM\cap B)\le d\Bigl(1+\frac{2}{\varepsilon}\Bigr)^{2m},
  \qquad\text{i.e.}\qquad
  \log\cN_\varepsilon(\cM\cap B)\le\log d+2m\log\Bigl(1+\frac{2}{\varepsilon}\Bigr),
\]
and the centres can be taken in $\cM\cap B$. The same bound holds for
$T(\cM)\cap B_W$, for every linear map $T\colon V\to W$ to a Euclidean space $W$
with unit ball $B_W$. In particular, it holds for the $L^2(\mu)$-seminorm, for
any probability measure $\mu$ on the inputs with $V\subseteq L^2(\mu)$, e.g.\ an
empirical measure.
\end{theorem}

For learning, let $V$ consist of functions on an input set $\cX$, let
\[
  \cF:=\{f\in\cM:\ |f(x)|\le 1\ \text{for all } x\in\cX\},
\]
let $(x_1,y_1),\dots,(x_n,y_n)$ be i.i.d.\ from a distribution $\pi$ on
$\cX\times[-1,1]$, and write $\cL_\pi(f)=\mathbb{E}_\pi(f(x)-y)^2$ and
$\cL_\cD(f)=\frac1n\sum_i(f(x_i)-y_i)^2$.

\begin{corollary}\label{cor:sample}
There is an absolute constant $C$ such that, with probability at least
$1-\delta$,
\[
  \sup_{f\in\cF}\bigl|\cL_\pi(f)-\cL_\cD(f)\bigr|
  \le C\sqrt{\frac{m+\log d+\log(2/\delta)}{n}}.
\]
In particular, $n=O\bigl(\varepsilon^{-2}(m+\log d+\log(1/\delta))\bigr)$
samples suffice for empirical risk minimisation over $\cF$ to be
$\varepsilon$-optimal in~$\cF$.
\end{corollary}

\section{Background}

Bounding the metric entropy of a set by its algebraic or integral-geometric
complexity has a long history. After Kolmogorov and Tikhomirov made
$\varepsilon$-entropy a systematic tool~\cite{KT59}, Vitushkin~\cite{Vi55}
introduced \emph{multidimensional variations} (integrals over affine sections of
the number of connected components) and showed that they bound covering numbers;
see also Ivanov~\cite{Iv75}. For semi-algebraic sets, the variations are
controlled by the Petrovskii--Oleinik--Milnor--Thom bounds on Betti
numbers~\cite{PO49,Mi64,Th65}. This yields
$\log\cN_\varepsilon\le m\log(1/\varepsilon)+C$, with $C$ depending on the
ambient dimension and on the degrees of the defining equations.
Yomdin~\cite{Yo87} used semi-algebraic entropy estimates to compare volume
growth and entropy of smooth maps, and Gromov's Bourbaki
exposition~\cite{Gr87} distilled the resulting ``algebraic lemma''; the
monograph of Yomdin and Comte~\cite{YC04} is the standard reference. A second
route goes through volumes of tubes, from Weyl's formula~\cite{We39} to
effective bounds for real algebraic sets by Wongkew~\cite{Wo93},
B\"urgisser, Cucker and Lotz~\cite{BCL08}, Lotz~\cite{Lo15}, and Basu and
Lerario~\cite{BL23}. Zhang and Kileel~\cite{ZK23} give a direct slicing argument
with improved bounds. In learning theory, covering numbers control sample
complexity~\cite{CS02}, and such bounds were brought to polynomial networks
in~\cite{KTB19,MSMTK25}.

On the complex side, Lelong~\cite{Le57} showed that integration over a complex
analytic set defines a closed positive current whose density ratio is monotone,
and Thie~\cite{Th67} identified the limiting density with the multiplicity.
Gromov~\cite{Gr03} (preprint 1977) combined Lelong's lower volume bound with a
packing argument to bound the topological entropy of holomorphic maps; the
argument is now standard in complex dynamics~\cite{DS10}. B\"urgisser, Cucker
and Lotz~\cite{BCL06} bounded volumes of tubes around complex projective
varieties in terms of their degree. The proof below transplants Gromov's packing
argument to the complexification of a real semi-algebraic set, where the total
volume is controlled by the degree through Stoll's theorem~\cite{St64}.

\section{Proofs}

\begin{lemma}[Lelong, Thie, Stoll]\label{lem:lelong}
Let $Y\subseteq\C^N$ be an irreducible algebraic variety of dimension $k$ and
degree $e$, and let $\omega_{2k}=\pi^k/k!$ be the volume of the unit ball of
$\C^k=\R^{2k}$. For $a\in\C^N$, the function
\[
  \Theta_a(r)=\frac{\cH^{2k}\bigl(Y\cap B(a,r)\bigr)}{\omega_{2k}\,r^{2k}},
  \qquad r>0,
\]
is nondecreasing, with $\lim_{r\to0}\Theta_a(r)=\mult_a Y\ge1$ if $a\in Y$, and
$\lim_{r\to\infty}\Theta_a(r)=e$. Consequently, for all $r>0$,
\begin{align}
  \cH^{2k}\bigl(Y\cap B(a,r)\bigr) &\ge \omega_{2k}\,r^{2k}
    \qquad (a\in Y), \label{eq:lower}\\
  \cH^{2k}\bigl(Y\cap B(a,r)\bigr) &\le e\,\omega_{2k}\,r^{2k}
    \qquad (a\in\C^N). \label{eq:upper}
\end{align}
\end{lemma}

\begin{proof}[Sketch]
By the Lelong--Jensen formula, for $0<s<r$,
\[
  \Theta_a(r)-\Theta_a(s)=\int_{Y\cap\{s<|z-a|<r\}}
  \bigl(dd^c\log|z-a|^2\bigr)^k\;\ge\;0,
\]
where $dd^c\log|z-a|^2$ is the pullback of the Fubini--Study form (of mass $1$
on lines) under $z\mapsto[z-a]\in\PP^{N-1}$. Hence $\Theta_a$ is nondecreasing,
and $\Theta_a(0^+)=\mult_aY$ by Thie's theorem~\cite{Th67}. As $s\to0$ and
$r\to\infty$, Crofton's formula on $\PP^{N-1}$ identifies the integral with the
number of points of $Y\cap L$ other than $a$, for a generic linear subspace
$L\ni a$ of codimension $k$. This number is $e-\mult_aY$, so
$\Theta_a(\infty)=e$, which is Stoll's theorem~\cite{St64}. See~\cite{Ch89,De12}
for details.
\end{proof}

\begin{proof}[Proof of Theorem~\ref{thm:main}]
Extend the inner product to a Hermitian one on $V_\C$, so that
$\|u+iv\|^2=\|u\|^2+\|v\|^2$ and $V\subset V_\C$ isometrically. In orthonormal
coordinates $V_\C=\C^N$, and dimension and degree are unchanged. Write
$m_j=\dim X_j\le m$ and $d_j=\deg X_j$, and partition
$\cM\cap B=S_1\sqcup\dots\sqcup S_s$ with $S_j\subseteq X_j$.

Let $P\subseteq S_j$ be finite and $\varepsilon$-separated, i.e.\
$\|f-g\|>\varepsilon$ for distinct $f,g\in P$. The complex balls
$B(f,\varepsilon/2)$, $f\in P$, are pairwise disjoint, centred on $X_j$, and
contained in $B(0,1+\varepsilon/2)$. By~\eqref{eq:lower} and~\eqref{eq:upper},
\[
  |P|\,\omega_{2m_j}\Bigl(\frac{\varepsilon}{2}\Bigr)^{2m_j}
  \le\sum_{f\in P}\cH^{2m_j}\bigl(X_j\cap B(f,\tfrac{\varepsilon}{2})\bigr)
  \le\cH^{2m_j}\bigl(X_j\cap B(0,1+\tfrac{\varepsilon}{2})\bigr)
  \le d_j\,\omega_{2m_j}\Bigl(1+\frac{\varepsilon}{2}\Bigr)^{2m_j},
\]
so $|P|\le d_j(1+2/\varepsilon)^{2m_j}$. Hence a maximal $\varepsilon$-separated
subset of $S_j$ exists, is finite, and is an $\varepsilon$-net of $S_j$ with
centres in $S_j$. Summing over $j$, using $m_j\le m$ and $\sum_j d_j=d$, gives
the bound.

For the last claim, extend $T$ $\C$-linearly. The complex Zariski closure of
$T(\cM)$ is $\overline{T(X)}$, whose dimension is at most $m$ and whose degree
is at most $d$, since linear maps do not increase degree~\cite{He83}. Apply the
first part to $T(\cM)\subseteq W$.
\end{proof}

\begin{proof}[Proof of Corollary~\ref{cor:sample}]
Fix $S=(x_1,\dots,x_n)$ and let $T_S\colon V\to\R^n$, $f\mapsto(f(x_i))_i$,
which is linear. With $\|u\|_S=(\frac1n\sum_i u_i^2)^{1/2}$ we have
$\cF|_S\subseteq T_S(\cM)\cap\{\|u\|_S\le1\}$, so Theorem~\ref{thm:main} gives,
uniformly in $S$ and $n$,
\[
  \log\cN_t\bigl(\cF|_S,\|\cdot\|_S\bigr)\le\log d+2m\log(1+4/t)
  \qquad(t>0),
\]
with centres in $\cF|_S$ (the $4$ accounts for moving the centres into
$\cF|_S$). Dudley's entropy integral~\cite{Du67} bounds the empirical
Rademacher complexity:
\[
  \widehat{\mathfrak R}_S(\cF)\lesssim\frac{1}{\sqrt n}\int_0^1
  \sqrt{\log d+2m\log(1+4/t)}\,dt\lesssim\sqrt{\frac{m+\log d}{n}}.
\]
The loss $(\hat y,y)\mapsto(\hat y-y)^2$ takes values in $[0,4]$ and is
$4$-Lipschitz in $\hat y\in[-1,1]$, so the contraction lemma and the standard
symmetrisation and bounded-differences argument~\cite{BM02} give the claim.
Uniform deviation at most $\varepsilon/2$ makes empirical risk minimisation
$\varepsilon$-optimal.
\end{proof}

\section{Comparison with Kileel et al.}

The neuroalgebraic geometry literature quotes~\cite[eq.~(4)]{MSMTK25},
attributing it to~\cite{KTB19},
\[
  \log\cN_\varepsilon(\cM)=O\bigl(m\log(d/\varepsilon)+C\bigr),
\]
with $C$ depending on $N=\dim V$. Zhang and Kileel~\cite{ZK23} prove bounds in
terms of the degree $D$ of the polynomials defining or parametrising the set.
Table~\ref{tab:comparison} lists them in the regime $m\ll N$, as summarised
in~\cite[Table~1]{ZK23} (in our notation). Its last two rows use two
Bézout-type facts.


\begin{itemize}
\item[(a)] If $\varphi\colon\C^q\to\C^N$ has coordinates of degree $\le D$ and
  $Y=\overline{\varphi(\C^q)}$ has dimension $m$, then $\deg Y\le D^m$. Indeed,
  for a generic affine subspace $L$ of codimension $m$, the $\deg Y$ points of
  $Y\cap L$ lie in $\varphi(\C^q)$. Their fibres partition $\varphi^{-1}(L)$
  into $\deg Y$ nonempty unions of irreducible components, and since
  $\varphi^{-1}(L)$ is cut out by $m$ equations of degree $\le D$, the Bézout
  inequality~\cite{He83} bounds the number of components by $D^m$.
\item[(b)] The zero set in $\C^N$ of polynomials of degree $\le D$ has degree
  at most $D^N$~\cite{He83}. The proof of Theorem~\ref{thm:main} only uses
  $\cM\subseteq X$, so it applies with this zero set in place of $X$.
\end{itemize}

\begin{table}[ht]
\centering\small
\renewcommand{\arraystretch}{1.3}
\begin{tabular}{@{}lll@{}}
\toprule
Source & Setting & Bound on $\log\cN_\varepsilon(\,\cdot\,\cap B)$\\
\midrule
\cite{KTB19} via \cite{MSMTK25} & variety of dimension $m$, degree $d$
  & $O\bigl(m\log\frac{d}{\varepsilon}+C(N)\bigr)$\\
\cite{ZK23} & image of $\varphi\colon\R^q\to\R^N$, $\deg\varphi\le D$
  & $q\log\frac1\varepsilon+O(q\log D+q\log N)$\\
\cite{ZK23} & real variety of dimension $m$, equations of degree $\le D$
  & $m\log\frac1\varepsilon+O(N\log D+m\log N)$\\
\midrule
Thm.~\ref{thm:main} & semi-algebraic, dimension $m$, $d=\deg\overline{\cM}$
  & $\log d+2m\log(1+\frac{2}{\varepsilon})$\\
Thm.~\ref{thm:main} + (a) & image as in row 2, $m=\dim\varphi(\R^q)\le q$
  & $m\log D+2m\log(1+\frac{2}{\varepsilon})$\\
Thm.~\ref{thm:main} + (b) & variety as in row 3, complex zero set of dim.\ $m$
  & $N\log D+2m\log(1+\frac{2}{\varepsilon})$\\
\bottomrule
\end{tabular}
\caption{Covering-number bounds for sets $U\cap B\subseteq\R^N$.}
\label{tab:comparison}
\end{table}

Compared with these results, Theorem~\ref{thm:main}:
\begin{itemize}
\item is stated in terms of the degree $d$ of the variety itself, which enters
  only as $\log d$. This removes the cross term $m\log d$ of~\cite{KTB19}, and
  by (a) and (b) it recovers the degree terms of~\cite{ZK23}. Moreover,
  $\log d$ can be much smaller than these Bézout bounds;
\item has no dependence on the ambient dimension $N$;
\item for neuromanifolds, replaces the number of parameters $q$ by
  $m=\dim\cM$, which is smaller whenever the parametrisation has
  positive-dimensional fibres (e.g.\ rescaling symmetries);
\item pays a factor $2$ in front of $m\log(1/\varepsilon)$, whereas the real
  bounds have the sharp leading term as $\varepsilon\to0$.
\end{itemize}
The decoupling of $m$ and $\log d$ matters for deep polynomial networks, whose
dimension grows linearly with depth while their degree grows
(super-)exponentially~\cite{SMK25}.

\section{Remarks}

\paragraph{Sharpness.} The form $\log d+\Theta(m)\log(1/\varepsilon)$ cannot be
improved. The $d$ parallel affine $m$-planes through
$\tfrac12e_1,\dots,\tfrac12e_d$ in $\R^{d+m}$ form a variety of degree $d$ that
needs at least $d\,(c/\varepsilon)^m$ balls for $\varepsilon<1/4$. Irreducible
examples exist too: the trigonometric moment curve
$\gamma(t)=k^{-1/2}(\cos t,\sin t,\dots,\cos kt,\sin kt)$ is a rational normal
curve of degree $d=2k$ on the unit sphere, and the points $\gamma(2\pi j/k)$,
$0\le j<k$, are orthonormal. Hence $\cN_\varepsilon\ge d/2$ for
$\varepsilon<1/\sqrt2$.

\paragraph{Why complexify.} Over $\R$ there is no analogue
of~\eqref{eq:lower}: $\{x^2+y^2=0\}$ is a point and $\{x^2+y^2=\eta\}$ is an
arbitrarily small oval. Real proofs therefore count connected components of
slices, and the Petrovskii--Oleinik--Milnor--Thom bounds used for this are the
source of the ambient and cross terms in Table~\ref{tab:comparison}. Complex
varieties are calibrated, hence minimal, so monotonicity supplies the lower bound
for free; the price is the factor $2$.

\paragraph{Sup-norm covers.} The Cucker--Smale bound~\cite{CS02}, as used
in~\cite[Thm.~B.1]{MSMTK25}, needs covers in
$\|f\|_\infty=\sup_{\cX}|f|$, assumed to be a norm on $V$. Take John's
ellipsoid~\cite{Jo48} $E$ with $E\subseteq\{\|f\|_\infty\le1\}\subseteq\sqrt N E$.
Theorem~\ref{thm:main} for the norm with unit ball $E$ gives
$\log\cN^{\infty}_\varepsilon(\cF)\le\log d+2m\log(1+4\sqrt N/\varepsilon)$,
with centres in $\cF$.

\begin{thebibliography}{99}\footnotesize\setlength{\itemsep}{1pt}

\bibitem{BM02} P.~L. Bartlett and S.~Mendelson, Rademacher and Gaussian
complexities: risk bounds and structural results, \emph{J. Mach. Learn. Res.}
3 (2002), 463--482.

\bibitem{BL23} S.~Basu and A.~Lerario, Hausdorff approximations and volume of
tubes of singular algebraic sets, \emph{Math. Ann.} 387 (2023), 79--109.

\bibitem{BCR98} J.~Bochnak, M.~Coste and M.-F. Roy, \emph{Real Algebraic
Geometry}, Springer, 1998.

\bibitem{BCL06} P.~B\"urgisser, F.~Cucker and M.~Lotz, Smoothed analysis of
complex conic condition numbers, \emph{J. Math. Pures Appl.} 86 (2006),
293--309.

\bibitem{BCL08} P.~B\"urgisser, F.~Cucker and M.~Lotz, The probability that a
slightly perturbed numerical analysis problem is difficult, \emph{Math. Comp.}
77 (2008), 1559--1583.

\bibitem{Ch89} E.~M. Chirka, \emph{Complex Analytic Sets}, Kluwer, 1989.

\bibitem{CS02} F.~Cucker and S.~Smale, On the mathematical foundations of
learning, \emph{Bull. Amer. Math. Soc.} 39 (2002), 1--49.

\bibitem{De12} J.-P. Demailly, \emph{Complex Analytic and Differential
Geometry}, book manuscript, Institut Fourier, 2012, Ch.~III.

\bibitem{DS10} T.-C. Dinh and N.~Sibony, Dynamics in several complex variables:
endomorphisms of projective spaces and polynomial-like mappings, in
\emph{Holomorphic Dynamical Systems}, Lecture Notes in Math. 1998, Springer,
2010; arXiv:0810.0811.

\bibitem{Du67} R.~M. Dudley, The sizes of compact subsets of Hilbert space and
continuity of Gaussian processes, \emph{J. Funct. Anal.} 1 (1967), 290--330.

\bibitem{Gr87} M.~Gromov, Entropy, homology and semialgebraic geometry (after
Y.~Yomdin), S\'eminaire Bourbaki, Exp.~663, \emph{Ast\'erisque} 145--146
(1987), 225--240.

\bibitem{Gr03} M.~Gromov, On the entropy of holomorphic maps, \emph{Enseign.
Math.} 49 (2003), 217--235 (preprint 1977).

\bibitem{He83} J.~Heintz, Definability and fast quantifier elimination in
algebraically closed fields, \emph{Theoret. Comput. Sci.} 24 (1983), 239--277.

\bibitem{Iv75} L.~D. Ivanov, \emph{Variations of Sets and Functions} (in
Russian), Nauka, Moscow, 1975.

\bibitem{Jo48} F.~John, Extremum problems with inequalities as subsidiary
conditions, in \emph{Studies and Essays Presented to R.~Courant}, Interscience,
1948, 187--204.

\bibitem{KTB19} J.~Kileel, M.~Trager and J.~Bruna, On the expressive power of
deep polynomial neural networks, \emph{NeurIPS} 32 (2019).

\bibitem{KT59} A.~N. Kolmogorov and V.~M. Tikhomirov, $\varepsilon$-entropy and
$\varepsilon$-capacity of sets in function spaces, \emph{Uspekhi Mat. Nauk}
14 (1959), no.~2, 3--86.

\bibitem{Le57} P.~Lelong, Int\'egration sur un ensemble analytique complexe,
\emph{Bull. Soc. Math. France} 85 (1957), 239--262.

\bibitem{Lo15} M.~Lotz, On the volume of tubular neighborhoods of real
algebraic varieties, \emph{Proc. Amer. Math. Soc.} 143 (2015), 1875--1889.

\bibitem{MSMTK25} G.~L. Marchetti, V.~Shahverdi, S.~Mereta, M.~Trager and
K.~Kohn, Algebra unveils deep learning -- an invitation to neuroalgebraic
geometry, \emph{ICML} 2025; arXiv:2501.18915.

\bibitem{Mi64} J.~Milnor, On the Betti numbers of real varieties, \emph{Proc.
Amer. Math. Soc.} 15 (1964), 275--280.

\bibitem{PO49} I.~G. Petrovskii and O.~A. Oleinik, On the topology of real
algebraic surfaces, \emph{Izv. Akad. Nauk SSSR Ser. Mat.} 13 (1949), 389--402.

\bibitem{SMK25} V.~Shahverdi, G.~L. Marchetti and K.~Kohn, On the geometry and
optimization of polynomial convolutional networks, \emph{AISTATS} 2025.

\bibitem{St64} W.~Stoll, The growth of the area of a transcendental analytic
set I, II, \emph{Math. Ann.} 156 (1964).

\bibitem{Th65} R.~Thom, Sur l'homologie des vari\'et\'es alg\'ebriques
r\'eelles, in \emph{Differential and Combinatorial Topology}, Princeton Univ.
Press, 1965.

\bibitem{Th67} P.~R. Thie, The Lelong number of a point of a complex analytic
set, \emph{Math. Ann.} 172 (1967), 269--312.

\bibitem{Vi55} A.~G. Vitushkin, \emph{On Multidimensional Variations} (in
Russian), Gostekhizdat, Moscow, 1955.

\bibitem{We39} H.~Weyl, On the volume of tubes, \emph{Amer. J. Math.} 61
(1939), 461--472.

\bibitem{Wo93} R.~Wongkew, Volumes of tubular neighbourhoods of real algebraic
varieties, \emph{Pacific J. Math.} 159 (1993), 177--184.

\bibitem{Yo87} Y.~Yomdin, Volume growth and entropy, \emph{Israel J. Math.} 57
(1987), 285--300; $C^k$-resolution of semialgebraic mappings, \emph{ibid.},
301--317.

\bibitem{YC04} Y.~Yomdin and G.~Comte, \emph{Tame Geometry with Application in
Smooth Analysis}, Lecture Notes in Math. 1834, Springer, 2004.

\bibitem{ZK23} Y.~Zhang and J.~Kileel, Covering number of real algebraic
varieties and beyond: improved bounds and applications, arXiv:2311.05116, 2023.

\end{thebibliography}

\end{document}
